\section{Conclusion}
\label{sec:conclusion}

This study examined whether a tabular foundation model and its explanations are justified for
capacity-limited student outreach. \tabicl ranked students well (ROC-AUC 0.957), but its differences
from XGBoost, logistic regression and an interpretable EBM lay within paired-bootstrap uncertainty;
its additional complexity is therefore unsupported on this task. Its SHAP explanations were
faithful under deletion but sensitive to the random seed and comparison group, and mostly restated
weak first-semester results. The human study found no evidence that charts improved predictions of
the flag (difference $-0.016$), while confidence was slightly higher and one unflagged case had
lower accuracy with charts.

Overall, the evidence favours a transparent model for further validation to \emph{order}, not
determine, a contact list. A next-stage design should combine a zero-pass check with transparent
ranking, display the cohort cut-off and explanation uncertainty, and retain adviser review.
Evidence from one institution and proxy users does not establish deployment validity.
Current-cohort evaluation with support staff is required; the system must never deny support.
